\section{Week 13: Hierarchical Models}

\begin{center}\rule{0.5\linewidth}{0.5pt}\end{center}

\subsection{Key terms}

Important terms of this lesson, with their Persian equivalents (repeated
even if they appeared in earlier weeks).

\begin{longtable}[]{@{}
  >{\raggedright\arraybackslash}p{(\linewidth - 4\tabcolsep) * \real{0.3333}}
  >{\raggedright\arraybackslash}p{(\linewidth - 4\tabcolsep) * \real{0.3333}}
  >{\raggedright\arraybackslash}p{(\linewidth - 4\tabcolsep) * \real{0.3333}}@{}}
\caption{\centering Key terms}\tabularnewline
\toprule\noalign{}
\begin{minipage}[b]{\linewidth}\raggedright
English
\end{minipage} & \begin{minipage}[b]{\linewidth}\raggedright
\RL{\CMfa فارسی}
\end{minipage} & \begin{minipage}[b]{\linewidth}\raggedright
\RL{\CMfa توضیح کوتاه}
\end{minipage} \\
\midrule\noalign{}
\endfirsthead
\toprule\noalign{}
\begin{minipage}[b]{\linewidth}\raggedright
English
\end{minipage} & \begin{minipage}[b]{\linewidth}\raggedright
\RL{\CMfa فارسی}
\end{minipage} & \begin{minipage}[b]{\linewidth}\raggedright
\RL{\CMfa توضیح کوتاه}
\end{minipage} \\
\midrule\noalign{}
\endhead
\bottomrule\noalign{}
\endlastfoot
hierarchical model & \RL{\CMfa مدل سلسله‌مراتبی (چندسطحی)} & \RL{\CMfa پارامترهای
فردی از یک توزیع گروهی می‌آیند} \\
parameter & \RL{\CMfa پارامتر} & \RL{\CMfa کمیت قابل تنظیم مدل که بین افراد یا شرایط
فرق می‌کند} \\
estimate / estimator & \RL{\CMfa برآورد} / \RL{\CMfa برآوردگر} & \RL{\CMfa مقدار به‌دست‌آمده
برای پارامتر} / \RL{\CMfa روش به دست آوردن آن} \\
trial & \RL{\CMfa کوشش} (trial) & \RL{\CMfa یک بار انجام تکلیف توسط شرکت‌کننده} \\
prior & \RL{\CMfa توزیع پیشین} (prior) & \RL{\CMfa باور درباره‌ی پارامتر پیش از دیدن
داده} \\
likelihood & \RL{\CMfa درست‌نمایی} & \RL{\CMfa احتمال داده به‌عنوان تابعی از
پارامترها} \\
prior predictive check & \RL{\CMfa بررسی پیش‌بینی پیشین} & \RL{\CMfa شبیه‌سازی داده از
پیشین برای سنجش معقول بودن آن} \\
posterior & \RL{\CMfa توزیع پسین} (posterior) & \RL{\CMfa باور درباره‌ی پارامتر بعد
از دیدن داده} \\
posterior predictive check (PPC) & \RL{\CMfa بررسی پیش‌بینی پسین} & \RL{\CMfa مقایسه‌ی
داده‌ی شبیه‌سازی‌شده از پسین با داده‌ی واقعی} \\
HMC / NUTS & Hamiltonian Monte Carlo / NUTS & MCMC \RL{\CMfa مبتنی بر
گرادیان} \\
divergent transition & divergence \RL{\CMfa (گذار واگرا)} & \RL{\CMfa هشدار} NUTS
\RL{\CMfa برای ناحیه‌ای که درست کاوش نشده} \\
complete / no / partial pooling & \RL{\CMfa ادغام کامل} / \RL{\CMfa بدون ادغام} /
\RL{\CMfa ادغام جزئی} (pooling) & \RL{\CMfa یک پارامتر برای همه} / \RL{\CMfa جدا برای هر
فرد} / \RL{\CMfa سلسله‌مراتبی} \\
shrinkage & \RL{\CMfa انقباض به سمت میانگین} (shrinkage) & \RL{\CMfa کشیده شدن
برآوردهای پرنویز به سمت میانگین گروه} \\
centered / non-centered parameterization & \RL{\CMfa پارامتری‌سازی} centered /
non-centered & \RL{\CMfa دو شکل نوشتن مدل سلسله‌مراتبی} \\
group-level parameter & \RL{\CMfa پارامتر سطح گروه} & \RL{\CMfa میانگین و پراکندگی
پارامترها در جمعیت} \\
fit / fitting & \RL{\CMfa برازش} & \RL{\CMfa تنظیم پارامترها تا مدل با داده جور
شود} \\
hit & \RL{\CMfa اصابت} (hit) & \RL{\CMfa پاسخ «بله» وقتی علامت وجود دارد} \\
sensitivity (d′) & \RL{\CMfa حساسیت} & \RL{\CMfa فاصله‌ی توزیع علامت و نویز} \\
gradient & \RL{\CMfa گرادیان} & \RL{\CMfa جهت و شیب تغییر تابع} \\
simulation & \RL{\CMfa شبیه‌سازی} & \RL{\CMfa تولید داده‌ی مصنوعی با اجرای مدل} \\
reparameterization & \RL{\CMfa بازپارامتری‌سازی} & \RL{\CMfa نوشتن پارامتر در مقیاس
دیگر، مثلاً} \ensuremath{\lambda} = e\^{}\ensuremath{\eta} \\
standard error (SE) & \RL{\CMfa خطای استاندارد} & \RL{\CMfa پراکندگی برآورد در
تکرارهای فرضی آزمایش} \\
benchmark & benchmark \RL{\CMfa (مجموعه‌آزمون استاندارد)} & \RL{\CMfa مجموعه سؤال
برای مقایسه‌ی مدل‌های هوش مصنوعی} \\
\end{longtable}

\begin{center}\rule{0.5\linewidth}{0.5pt}\end{center}

\subsection{Learning outcomes}

By the end of this week, students will be able to:

\begin{enumerate}
\def\labelenumi{\arabic{enumi}.}
\tightlist
\item
  \textbf{Explain} conceptually how Hamiltonian Monte Carlo / NUTS
  differs from random-walk Metropolis, and what a divergence signals.
\item
  \textbf{Formulate} a hierarchical (multilevel) model and
  \textbf{explain} partial pooling and shrinkage, contrasting it with
  the complete- and no-pooling approaches of Week 5.
\item
  \textbf{Formulate and evaluate} a non-centered parameterization and
  \textbf{justify} its use.
\item
  \textbf{Interpret} group-level and individual-level posteriors and
  \textbf{report} them according to the Week 12 checklist.
\end{enumerate}

\begin{center}\rule{0.5\linewidth}{0.5pt}\end{center}

\subsection{Motivation \& big picture}

In Week 5 we faced a dilemma. Fitting the \textbf{average} of all
participants is stable but can distort the model (aggregation artifacts)
and ignores individual differences. Fitting \textbf{each participant
separately} respects differences, but with few trials per person the
estimates are noisy. Some become extreme or infinite (a participant with
50/50 hits has \(d' = \infty\) without a correction). The second-stage
statistics on those estimates also ignore that some are much less
certain than others.

\textbf{Hierarchical models} dissolve the dilemma. They assume that each
participant has their own parameter, \emph{and} that these individual
parameters come from a group-level distribution whose mean and spread
are also estimated. Each individual estimate then borrows strength from
the others: noisy, extreme estimates are pulled (``shrunk'') toward the
group mean, by exactly as much as their uncertainty warrants. The
group-level parameters answer the population question (``what is the
typical sensitivity, and how much do people differ?'') with all sources
of uncertainty included.

Writing samplers for such models by hand would be a research project.
The hand-built samplers of Weeks 11--12 are what make this black box
transparent to you.

\paragraph{Where is this used?}

\begin{itemize}
\tightlist
\item
  \textbf{Most modern cognitive modeling studies with many participants}
  fit hierarchical models: hierarchical drift-diffusion models, hierarchical reinforcement-learning models in computational
  psychiatry, hierarchical SDT in memory research.
\item
  \textbf{Clinical and developmental studies with short sessions}:
  children and patients cannot do hundreds of trials, so individual
  estimates are noisy; partial pooling makes them usable.
\item
  \textbf{Education and social science}: students within classes within
  schools (multilevel models) are the textbook example of hierarchical
  structure.
\item
  \textbf{Sports analytics and medicine}: estimating a player's skill or
  a hospital's success rate from few observations by borrowing strength
  from the group is the same idea.
\end{itemize}

\begin{center}\rule{0.5\linewidth}{0.5pt}\end{center}

\Needspace{11\baselineskip}

\subsection{Lecture notes}

\paragraph{In plain words}

When each person gives us only a little data, their individual estimates
are noisy, and some come out absurd (a perfect score gives
\(d' = \infty\)). A \textbf{hierarchical model} says: each person has
their own parameter, \emph{but} people are similar to each other, so
their parameters come from a common group distribution, whose mean and
spread we also estimate. The effect is sensible compromise: a person
with little or extreme data is pulled toward the group average
(\textbf{shrinkage}), and a person with plenty of data keeps an estimate
close to their own. How strong the pull is depends on how different
people turn out to be, which is learned from the data. The last part of
the week is a step-by-step \textbf{workflow} for doing this safely.

\subsubsection{2. Hamiltonian Monte Carlo and NUTS (conceptually)}

Random-walk Metropolis proposes blind steps and wastes most of them in
high dimensions or along correlated ridges (Weeks 11--12).
\textbf{Hamiltonian Monte Carlo} (HMC) treats the negative log posterior
as a landscape and the current parameter vector as a frictionless ball.
Give the ball a random push (momentum), let it roll for a while
following the landscape's gradient, and propose where it ends up.
Proposals can be far away yet still in high-probability regions, so
acceptance is high and autocorrelation low. \textbf{NUTS} (the No-U-Turn
Sampler) tunes the rolling time automatically.

\textbf{Divergences.} The rolling is simulated numerically in small
steps. Where the posterior's curvature changes abruptly (typically in
``funnels'', below), the simulation becomes inaccurate and NUTS flags a
\textbf{divergent transition}. Even a few divergences mean the posterior
is not being explored reliably there. Do not ignore them;
reparameterize.

\begin{quote}
\textbf{Advanced:} HMC simulates Hamiltonian dynamics with energy
\(H(\boldsymbol\theta, \mathbf{p}) = -\ln\tilde p(\boldsymbol\theta) + \tfrac12\mathbf{p}^\top M^{-1}\mathbf{p}\)
using the leapfrog integrator. Exact dynamics conserve \(H\), so
proposals would always be accepted. The leapfrog error makes a
Metropolis correction necessary, and a divergence is a trajectory whose
energy error explodes.
\end{quote}

\subsubsection{3. Hierarchical models}

\textbf{Structure (for SDT with \(J\) participants):}

\[
\begin{aligned}
&\text{Group level:} && \mu_d \sim \mathcal{N}(1, 1),\ \ \sigma_d \sim \mathcal{N}^{+}(0, 1),\ \ \mu_c \sim \mathcal{N}(0, 1),\ \ \sigma_c \sim \mathcal{N}^{+}(0, 1)\\
&\text{Individual level:} && d'_j \sim \mathcal{N}(\mu_d, \sigma_d),\ \ c_j \sim \mathcal{N}(\mu_c, \sigma_c), \quad j = 1, \dots, J\\
&\text{Data:} && \text{hits}_j \sim \mathrm{Binomial}\big(n_s, \Phi(d'_j/2 - c_j)\big),\ \ \text{FAs}_j \sim \mathrm{Binomial}\big(n_n, \Phi(-d'_j/2 - c_j)\big)
\end{aligned}
\]

\textbf{Three ways to pool information (cf.~Week 5):}

\begin{longtable}[]{@{}
  >{\raggedright\arraybackslash}p{(\linewidth - 4\tabcolsep) * \real{0.3333}}
  >{\raggedright\arraybackslash}p{(\linewidth - 4\tabcolsep) * \real{0.3333}}
  >{\raggedright\arraybackslash}p{(\linewidth - 4\tabcolsep) * \real{0.3333}}@{}}
\caption{\centering 3. Hierarchical models}\tabularnewline
\toprule\noalign{}
\begin{minipage}[b]{\linewidth}\raggedright
Approach
\end{minipage} & \begin{minipage}[b]{\linewidth}\raggedright
Model
\end{minipage} & \begin{minipage}[b]{\linewidth}\raggedright
What it assumes
\end{minipage} \\
\midrule\noalign{}
\endfirsthead
\toprule\noalign{}
\begin{minipage}[b]{\linewidth}\raggedright
Approach
\end{minipage} & \begin{minipage}[b]{\linewidth}\raggedright
Model
\end{minipage} & \begin{minipage}[b]{\linewidth}\raggedright
What it assumes
\end{minipage} \\
\midrule\noalign{}
\endhead
\bottomrule\noalign{}
\endlastfoot
Complete pooling & One \(d'\) and one \(c\) for everyone &
\(\sigma_d = \sigma_c = 0\): no individual differences \\
No pooling & Separate \(d'_j\), \(c_j\) with fixed, independent priors &
\(\sigma \to \infty\): participants tell us nothing about each other \\
\textbf{Partial pooling} & Hierarchical; \(\sigma_d\), \(\sigma_c\)
\emph{estimated} & The data decide how similar participants are \\
\end{longtable}

\textbf{Shrinkage.} For a normal approximation, the hierarchical
estimate of participant \(j\) is a precision-weighted average of their
own estimate \(\hat d_j\) (with standard error \(\mathrm{SE}_j\)) and
the group mean:

\[
\mathbb{E}[d'_j \mid \text{data}] \approx \frac{\hat d_j/\mathrm{SE}_j^2 + \mu_d/\sigma_d^2}{1/\mathrm{SE}_j^2 + 1/\sigma_d^2}.
\]

Participants with \textbf{noisier} data (larger \(\mathrm{SE}_j\)) are
shrunk \textbf{more}. If people really differ a lot (large
\(\sigma_d\)), shrinkage is weak; if they are similar, it is strong.
This is learned from the data.

\textbf{Why shrinkage is good, not a distortion:} across participants,
shrunken estimates are on average \emph{closer to the true values} than
the raw individual estimates. This is a famous result in statistics
(Stein's paradox). Extreme raw estimates are partly extreme by chance,
and shrinkage corrects for that regression to the mean.

\textbf{Group-level questions} are answered directly by the posteriors
of \(\mu_d\) (typical sensitivity) and \(\sigma_d\) (individual
differences), and group comparisons are made by adding group or
condition effects to the group level.

\paragraph{Reading the math}

\begin{itemize}
\tightlist
\item
  \textbf{Two levels.} Group level: \(\mu_d\) (typical sensitivity) and
  \(\sigma_d\) (how much people differ). Individual level:
  \(d'_j \sim \mathcal{N}(\mu_d, \sigma_d)\), person \(j\)'s sensitivity
  is drawn from the group distribution. Data level: the hits and false
  alarms of person \(j\) depend only on that person's \(d'_j\) and
  \(c_j\).
\item
  \(\mathcal{N}^{+}(0, 1)\): a ``half-normal'' distribution, a normal
  curve folded at 0 so only positive values remain (a standard deviation
  cannot be negative).
\item
  \textbf{The shrinkage formula} is a precision-weighted average (as in
  Week 12): the person's own estimate \(\hat d_j\) has precision
  \(1/\mathrm{SE}_j^2\), the group distribution has precision
  \(1/\sigma_d^2\). Example from the worked example:
  \(\mathrm{SE} = 0.40\) gives precision 6.25, the group gives 4, so the
  estimate is \((6.25 \times 2.4 + 4 \times 1.3)/10.25 = 1.97\): the raw
  2.4 moves 40\% of the way toward the group mean 1.3.
\item
  If people really differ a lot (large \(\sigma_d\)), the group
  precision is small and shrinkage is weak. If they are similar,
  shrinkage is strong.
\end{itemize}

\subsubsection{4. Centered vs.~non-centered parameterization}

Writing \(d'_j \sim \mathcal{N}(\mu_d, \sigma_d)\) directly (the
\textbf{centered} form) creates a ``funnel'' in the posterior: when
\(\sigma_d\) is small, all \(d'_j\) must be squeezed close together, and
when it is large they can spread out. NUTS struggles with this changing
geometry and produces divergences, especially with few data per
participant. The \textbf{non-centered} form samples standardized
deviations and builds \(d'_j\) deterministically:

\[
z_j \sim \mathcal{N}(0, 1), \qquad d'_j = \mu_d + \sigma_d z_j.
\]

It describes the same model but has a much friendlier geometry. Use it
by default when data per participant are sparse.

\paragraph{Reading the math}

\(z_j \sim \mathcal{N}(0, 1)\), \(d'_j = \mu_d + \sigma_d z_j\): instead
of sampling \(d'_j\) directly, sample a standardized ``how unusual is
this person'' score \(z_j\) (in units of the group SD) and rebuild
\(d'_j\) from it. Both versions describe exactly the same model. The
difference is only in what the sampler has to explore: in the centered
form, the allowed range of every \(d'_j\) shrinks to almost nothing when
\(\sigma_d\) is small (a funnel); in the non-centered form, the \(z_j\)
always live on the same scale, so the sampler's step size fits
everywhere.

\subsubsection{5. The Bayesian workflow}

\begin{enumerate}
\def\labelenumi{\arabic{enumi}.}
\tightlist
\item
  \textbf{Specify} the model and priors (with justification).
\item
  \textbf{Prior predictive check} (Week 9): does the prior imply
  plausible data?
\item
  \textbf{Fit} with several chains.
\item
  \textbf{Diagnose:} \(\hat R\), ESS, divergences, trace plots (Week
  12).
\item
  \textbf{Posterior predictive check} (Week 10): can the fitted model
  reproduce the data?
\item
  \textbf{Interpret and report.}
\end{enumerate}

\paragraph{For an engineer}

\begin{itemize}
\tightlist
\item
  \textbf{Hierarchical model = estimation with a population prior.} It
  is like estimating each user's channel gain from a few pilots while
  also learning the distribution of gains in the cell, and using it as
  the prior: an \textbf{empirical-Bayes / MMSE} estimator. Shrinkage
  toward the group mean is the same phenomenon as the
  \textbf{James--Stein} estimator, which beats separate ML estimates in
  total squared error when there are three or more units.
\item
  \textbf{HMC/NUTS} simulates Hamiltonian dynamics (a frictionless
  particle on the surface −log posterior) with a numerical integrator
  (leapfrog), as in physics simulations. A \textbf{divergence} is the
  integrator becoming unstable because the step size is too large for
  the local curvature, like a stiff ODE integrated with a too-large
  step.
\item
  \textbf{Non-centered parameterization} is a change of variables that
  normalizes the variables to unit scale, i.e.~it improves the
  conditioning of the problem, much like whitening or normalizing inputs
  before an adaptive algorithm.
\end{itemize}

\paragraph{For the instructor}

\begin{itemize}
\tightlist
\item
  \textbf{Introduce shrinkage with a sports or exam example.} A student
  who answered 3 of 3 questions correctly is probably not a ``100\%
  student''; a student with 95 of 100 probably is near 95\%. Ask the
  class how they would combine individual and class information: they
  invent partial pooling themselves.
\item
  \textbf{Show the shrinkage plot from lab Part D}: raw estimates,
  hierarchical estimates and the group mean on one figure, with arrows.
  It makes the idea stick.
\item
  \textbf{Do not teach HMC mathematics in depth.} One picture (a ball
  rolling on the negative log posterior, given a random push) is enough;
  the important practical skill is reading the diagnostics (divergences,
  \(\hat R\), ESS).
\item
  \textbf{Typical questions.}

  \begin{itemize}
  \tightlist
  \item
    \emph{``Isn't shrinkage biasing the results?''} Slightly, but it
    reduces the total error, so estimates are closer to the truth on
    average.
  \item
    \emph{``Is a hierarchical model always better?''} When individuals
    have enough data, results are similar to separate fits; the gain is
    largest with little data per person.
  \item
    \emph{``What does a divergence mean for my results?''} That part of
    the posterior was not explored reliably. Reparameterize and rerun
    until there are none.
  \end{itemize}
\end{itemize}

\begin{center}\rule{0.5\linewidth}{0.5pt}\end{center}

\subsection{Worked example}

\textbf{Shrinkage by hand.} The group posterior has
\(\mu_d \approx 1.3\) and \(\sigma_d \approx 0.5\) (group precision
\(1/0.5^2 = 4\)). Two participants both have a raw estimate
\(\hat d_j = 2.4\).

\begin{longtable}[]{@{}
  >{\raggedright\arraybackslash}p{(\linewidth - 8\tabcolsep) * \real{0.2000}}
  >{\raggedright\arraybackslash}p{(\linewidth - 8\tabcolsep) * \real{0.2000}}
  >{\raggedright\arraybackslash}p{(\linewidth - 8\tabcolsep) * \real{0.2000}}
  >{\raggedright\arraybackslash}p{(\linewidth - 8\tabcolsep) * \real{0.2000}}
  >{\raggedright\arraybackslash}p{(\linewidth - 8\tabcolsep) * \real{0.2000}}@{}}
\caption{\centering Worked example}\tabularnewline
\toprule\noalign{}
\begin{minipage}[b]{\linewidth}\raggedright
Participant
\end{minipage} & \begin{minipage}[b]{\linewidth}\raggedright
\(\mathrm{SE}_j\)
\end{minipage} & \begin{minipage}[b]{\linewidth}\raggedright
Individual precision \(1/\mathrm{SE}_j^2\)
\end{minipage} & \begin{minipage}[b]{\linewidth}\raggedright
Hierarchical estimate
\end{minipage} & \begin{minipage}[b]{\linewidth}\raggedright
Posterior SD
\end{minipage} \\
\midrule\noalign{}
\endfirsthead
\toprule\noalign{}
\begin{minipage}[b]{\linewidth}\raggedright
Participant
\end{minipage} & \begin{minipage}[b]{\linewidth}\raggedright
\(\mathrm{SE}_j\)
\end{minipage} & \begin{minipage}[b]{\linewidth}\raggedright
Individual precision \(1/\mathrm{SE}_j^2\)
\end{minipage} & \begin{minipage}[b]{\linewidth}\raggedright
Hierarchical estimate
\end{minipage} & \begin{minipage}[b]{\linewidth}\raggedright
Posterior SD
\end{minipage} \\
\midrule\noalign{}
\endhead
\bottomrule\noalign{}
\endlastfoot
P1 (few trials) & 0.40 & 6.25 &
\((6.25 \times 2.4 + 4 \times 1.3)/(6.25 + 4) = 20.2/10.25 = \mathbf{1.97}\)
& \(1/\sqrt{10.25} = 0.31\) \\
P2 (many trials) & 0.20 & 25 &
\((25 \times 2.4 + 4 \times 1.3)/(25 + 4) = 65.2/29 = \mathbf{2.25}\) &
\(1/\sqrt{29} = 0.19\) \\
\end{longtable}

\textbf{Interpretation.} Both participants \emph{look} equally
sensitive, but P1's estimate rests on less data, so the model trusts it
less and moves it 0.43 toward the group mean. P2 moves only 0.15. Both
posterior SDs are smaller than the raw SEs, because the group
distribution adds information. (In the full model, \(\mu_d\) and
\(\sigma_d\) are themselves uncertain and estimated jointly; this
calculation fixes them for illustration.)

\begin{center}\rule{0.5\linewidth}{0.5pt}\end{center}

\Needspace{11\baselineskip}

\subsection{Lab (2 hours)}

\subsubsection{Goals}

\begin{enumerate}
\def\labelenumi{\arabic{enumi}.}
\tightlist
\item
  Simulate a 12-participant experiment and compute individual
  (no-pooling) estimates.
\item
  Sample with NUTS; check \(\hat R\), ESS and divergences.
\item
  Visualize shrinkage and compare the estimates with the true values.
\item
  Run a posterior predictive check and interpret the group-level
  parameters.
\end{enumerate}

\textbf{Timing:} Part A 10 min · Part B 20 min · Part C 25 min · Part D
30 min · Part E 20 min · Part F 15 min.

\begin{quote}
Start Part C early.
\end{quote}

\subsubsection{Part B: prior predictive check}

The prior implies a broad range of hit and false-alarm rates, mostly
above the diagonal. About 30\% of simulated participants have a hit or
false-alarm count of exactly 0 or 50. These are mostly participants
simulated with very high sensitivity and a strong bias, which produce
zero false alarms. Try it and check whether the posterior changes.

\subsubsection{Part D: shrinkage and group-level inference (30 min)}

Which participants are shrunk most, and why? Find one whose raw estimate
is extreme and describe what the hierarchical model does with it. 2. Is
the hierarchical estimate closer to the truth on average? Is it closer
for \emph{every} participant? Why is that not required? 3. What do
\(\mu_d\) and \(\sigma_d\) tell you about this population? Compare
\(\sigma_d\) with the SD used to simulate the data (0.5) and explain the
difference (hint: 12 participants).

\subsubsection{Part F: report (15 min)}

Write a methods-and-results paragraph following the Week 12 checklist:
model and priors (in words and symbols), sampler and settings,
diagnostics (max \(\hat R\), min ESS, divergences), the PPC result, and
the group-level conclusions (\(\mu_d\), \(\sigma_d\), \(\mu_c\) with
95\% intervals; \(P(\mu_c > 0\mid\text{data})\)).

\subsubsection{Exercises (increasing difficulty)}

Compare its estimate of \(d'\) with \(\mu_d\) from the hierarchical
model, and its posterior predictive check with Part E. Which
participants does it fail to predict?

Compare divergences, ESS and run time. Then reduce the data to 10 + 10
trials per participant and compare again.

\textbf{Exercise 3 \ensuremath{\star}\ensuremath{\star} (unequal trial numbers).} Give half of the
participants 20 + 20 trials and the other half 200 + 200 (simulate new
data). Show that shrinkage depends on the number of trials, as in the
worked example.

Check the Pareto-\(k\) diagnostics and say what they mean for
leave-one-participant-out prediction.

\textbf{Exercise 5 \ensuremath{\star}\ensuremath{\star}\ensuremath{\star} (Advanced: a group effect).} Simulate two groups
of 12 participants that differ in mean sensitivity by 0.4. Extend the
model with a group effect \(\beta\) on \(\mu_d\)
(\(d'_j = \mu_d + \beta\cdot\text{group}_j + \sigma_d z_j\)), estimate
\(\beta\), and report \(P(\beta > 0\mid\text{data})\). Compare with a
t-test on the individual \(d'\) estimates.

\subsubsection{Lab check-off}

Show: (1) the diagnostics output (no divergences, \(\hat R\), ESS), (2)
the shrinkage plot with the RMSE comparison, and (3) the PPC plot with
your Part F paragraph.

\textbf{Project progress check (this week):} submit a one-page status
report (what is done, a first fit or simulation figure, open problems).
It is not graded; written feedback is returned within a week.

\begin{center}\rule{0.5\linewidth}{0.5pt}\end{center}

\subsection{Connection to AI \& LLMs}

\textbf{Partial pooling for AI benchmarks.} Benchmarks are split into
task categories, and some categories have only a handful of items.
Reporting raw accuracy per category gives noisy, sometimes extreme
numbers, just like individual SDT fits with few trials. A hierarchical
model treats category accuracies as draws from a common distribution and
shrinks small categories toward the overall mean.

The 8-item category (raw 0.88, true 0.63) is pulled down strongly, and
the shrunken estimates are about twice as accurate overall (RMSE 0.066
vs.~0.125).

Relate your answer to the shrinkage formula in this week's lecture
notes. Compare the posterior means with the empirical-Bayes estimates,
and check the diagnostics. - \ensuremath{\star}\ensuremath{\star} Discussion: a leaderboard highlights
``the category where Model X is weakest'', a category with 5 items. What
would you tell the authors?

\begin{center}\rule{0.5\linewidth}{0.5pt}\end{center}

\subsection{Common mistakes \& misconceptions}

\begin{enumerate}
\def\labelenumi{\arabic{enumi}.}
\tightlist
\item
  \textbf{``Shrinkage biases the results.''} It trades a little bias for
  a large reduction in variance and gives more accurate estimates on
  average. It is principled, not an ad-hoc correction.
\item
  \textbf{Ignoring divergences.} Even a handful indicate unreliable
  exploration.
\item
  \textbf{Centered parameterization with sparse data.} It invites
  funnels and divergences. Default to non-centered for hierarchical
  models with few trials per participant.
\item
  \textbf{Vague priors on group SDs.} Very wide priors on \(\sigma\) can
  produce implausible prior predictions and sampling problems.
  Half-normal priors on a sensible scale are usually better.
\item
  \textbf{Reporting only individual estimates.} The group-level
  parameters often answer the actual research question.
\end{enumerate}

\begin{center}\rule{0.5\linewidth}{0.5pt}\end{center}

\subsection{Check your understanding}

\textbf{1.} What are the three ways to combine data across participants,
and what does each assume about individual differences? \textbf{Answer}
Complete pooling (everyone identical, \(\sigma = 0\)), no pooling
(participants unrelated, so no information is shared), and partial
pooling/hierarchical (participants come from a common distribution whose
spread is estimated from the data).

\textbf{2.} Two participants have the same raw \(d'\) estimate. Why may
the hierarchical model shrink one of them much more? \textbf{Answer}
Shrinkage depends on the precision of each individual estimate. The
participant with fewer trials (larger SE) carries less information and
is pulled more strongly toward the group mean.

\textbf{3.} Why is the non-centered parameterization often preferred?
\textbf{Answer} The centered form creates a funnel-shaped posterior
(individual parameters must collapse together when the group SD is
small) that NUTS explores poorly, causing divergences. Sampling
standardized \(z_j\) and computing \(d'_j = \mu + \sigma z_j\) gives a
geometry that is much easier to explore.

\textbf{4.} What is a divergent transition, and what should you do about
it? \textbf{Answer} A NUTS trajectory whose numerical simulation broke
down, usually in a region of high curvature; it signals that part of the
posterior is not explored reliably.

\textbf{5.} Why does HMC usually mix much better than random-walk
Metropolis in high dimensions? \textbf{Answer} It uses gradient
information to make long, directed moves that follow the shape of the
posterior, so distant proposals still have a high acceptance
probability. Random-walk proposals become almost always rejected unless
they are tiny.

\textbf{6.} In a hierarchical SDT model, which parameter answers ``how
much do people differ in sensitivity?'', and which answers ``what is the
typical sensitivity?'' \textbf{Answer} \(\sigma_d\) (the group-level SD)
and \(\mu_d\) (the group-level mean), respectively.

\textbf{7.} List the steps of the Bayesian workflow. \textbf{Answer}
Specify the model and priors \ensuremath{\rightarrow} prior predictive check \ensuremath{\rightarrow} fit with several
chains \ensuremath{\rightarrow} diagnose (\(\hat R\), ESS, divergences, traces) \ensuremath{\rightarrow} posterior
predictive check \ensuremath{\rightarrow} compare models if needed \ensuremath{\rightarrow} interpret and report.

\begin{center}\rule{0.5\linewidth}{0.5pt}\end{center}

\subsection{Assignment}

\begin{itemize}
\tightlist
\item
  \textbf{A5 (MCMC from scratch)} is due at the start of the Week 14
  lecture.
\item
  \textbf{Project progress check} this week (see above).
\item
  No new graded assignment: the remaining weeks are reserved for the
  final project.
\end{itemize}

\begin{center}\rule{0.5\linewidth}{0.5pt}\end{center}

\subsection{Further reading}

\begin{itemize}
\tightlist
\item
  \textbf{Lee \& Wagenmakers (2014), chapter 11.} Signal detection
  theory, including the hierarchical SDT model you implemented today.
  \emph{(Main reading.)}
\item
  \textbf{Kruschke (2014), chapter 9.} Hierarchical models: partial
  pooling and shrinkage explained step by step.
\item
  \textbf{Farrell \& Lewandowsky (2018), chapter 9.} Multilevel or
  hierarchical modeling of cognitive data.
\end{itemize}

\emph{Preview for Week 14:} with the full toolkit in hand, we tour
models from different areas of cognitive science, starting with
categorization (the exemplar model GCM) and the speed--accuracy
machinery of the drift-diffusion model.
